\documentclass[11pt]{ctexart}
\usepackage[a4paper,margin=1in]{geometry}
\usepackage{longtable,booktabs}
\title{Geant4-Agent 回归测试报告（中文）}
\author{自动化评测流程}
\date{2026-03-06\_095451}
\begin{document}
\maketitle
\section{统一 Pass/Fail 总表}
\begin{longtable}{p{0.06\linewidth}p{0.08\linewidth}p{0.07\linewidth}p{0.08\linewidth}p{0.07\linewidth}p{0.07\linewidth}p{0.07\linewidth}p{0.07\linewidth}p{0.09\linewidth}p{0.06\linewidth}p{0.21\linewidth}}
\toprule
ID & 领域 & 风格 & 参数完整性 & 期望初始 & 实际初始 & 期望最终 & 实际最终 & 缺参(初->终) & 结果 & 失败原因 \\
\midrule
MTM01 & medical & fuzzy & missing & 失败 & 失败 & 通过 & 通过 & 13->0 & 通过 & - \\
MTM02 & aerospace & precise & missing & 失败 & 失败 & 通过 & 通过 & 2->0 & 通过 & - \\
MTM03 & industrial\_ndt & precise & missing & 失败 & 失败 & 通过 & 通过 & 4->0 & 失败 & expected\_final\_mismatch:geometry.structure \\
MTM04 & physics & precise & missing & 失败 & 失败 & 通过 & 通过 & 8->0 & 通过 & - \\
MTM05 & nuclear\_engineering & fuzzy & missing & 失败 & 失败 & 通过 & 失败 & 13->1 & 失败 & final\_complete\_mismatch, expected\_final\_mismatch:materials.selected\_materials.0 \\
MTM06 & security\_screening & fuzzy & missing & 失败 & 失败 & 通过 & 通过 & 4->0 & 失败 & expected\_final\_mismatch:geometry.structure \\
MTM07 & semiconductor & precise & missing & 失败 & 失败 & 通过 & 失败 & 13->1 & 失败 & final\_complete\_mismatch, expected\_final\_mismatch:materials.selected\_materials.0 \\
MTM08 & education & fuzzy & missing & 失败 & 失败 & 通过 & 通过 & 11->0 & 通过 & - \\
MTM09 & environmental\_radiation & fuzzy & missing & 失败 & 失败 & 通过 & 失败 & 13->1 & 失败 & final\_complete\_mismatch, expected\_final\_mismatch:materials.selected\_materials.0 \\
MTM10 & high\_energy\_physics & precise & missing & 失败 & 失败 & 通过 & 通过 & 10->0 & 通过 & - \\
\bottomrule
\end{longtable}
\section{基准对比}
\begin{itemize}
\item pass\_rate=0.500
\item initial\_alignment\_rate=1.000
\item final\_alignment\_rate=0.700
\item delta\_pass\_rate=-0.450
\item regressions=无
\item improvements=无
\end{itemize}
\appendix
\section{附录：完整对话与内部处理链（自然语言）}
\subsection{案例 MTM01（medical）}
\textbf{对话：}
\begin{itemize}
\item 用户：????????????????????????????????
\item 系统：您想构建哪种几何结构呢？比如一个简单的长方体，还是一个圆柱体？另外，长方体的X方向长度大概是多少？
\item 用户：???? 300 mm x 300 mm x 300 mm ??????? G4\_WATER?source.type=beam?particle=gamma?energy=6 MeV?position=(0,0,-400) mm?direction=+z?physics=QBBC?output=root?
\item 系统：配置已完成。
\end{itemize}
\textbf{内部链路（自然语言）：}
\begin{itemize}
\item 第 1 轮：动作=ask\_clarification，阶段=geometry，闭环状态=False。
\item 缺失字段：geometry.structure、geometry.params.module\_x、geometry.params.module\_y、geometry.params.module\_z、materials.volume\_material\_map、source.type、source.particle、source.energy、source.position、source.direction、physics.physics\_list、output.format、output.path。
\item LLM 使用=True；回退=无；后端=llm\_semantic\_frame+runtime\_semantic。
\item 第 2 轮：动作=finalize，阶段=materials，闭环状态=True。
\item 缺失字段：无。
\item LLM 使用=True；回退=无；后端=llm\_slot\_frame+runtime\_semantic。
\end{itemize}
\subsection{案例 MTM02（aerospace）}
\textbf{对话：}
\begin{itemize}
\item 用户：I need a first-pass proton irradiation setup for a silicon device; keep it simple for now.
\item 系统：Here's your current setup status:

**Geometry:** Updated with a simple structure and module dimensions.
**Materials:** Silicon is selected and assigned to your device volume.
**Source:** Configured for proton irradiation with a specified energy.
**Physics:** A suitable physics list has been selected.

**Next Steps:** To run the simulation, please provide:
*   The output format (e.g., ROOT, CSV)
*   The output file path

**Note:** Providing a path will create a new file. If a file already exists at that location, it will be overwritten.
\item 用户：Use a silicon box 50 mm x 50 mm x 5 mm, beam source proton 100 MeV from (0,0,-100) mm along +z, physics FTFP\_BERT, output json.
\item 系统：Your requested geometry changes would overwrite two existing values. Please confirm this update:

*   `module\_x`: 1000.0 → 50.0
*   `module\_y`: 1000.0 → 50.0
\item 用户：confirm
\item 系统：Configuration complete.
\end{itemize}
\textbf{内部链路（自然语言）：}
\begin{itemize}
\item 第 1 轮：动作=summarize\_progress，阶段=materials，闭环状态=False。
\item 缺失字段：output.format、output.path。
\item LLM 使用=True；回退=无；后端=llm\_slot\_frame+runtime\_semantic。
\item 第 2 轮：动作=confirm\_overwrite，阶段=source，闭环状态=False。
\item 缺失字段：无。
\item LLM 使用=True；回退=无；后端=llm\_slot\_frame+runtime\_semantic。
\item 第 3 轮：动作=finalize，阶段=physics，闭环状态=True。
\item 缺失字段：无。
\item LLM 使用=True；回退=无；后端=llm\_semantic\_frame+runtime\_semantic+confirmed\_pending\_overwrite。
\end{itemize}
\subsection{案例 MTM03（industrial\_ndt）}
\textbf{对话：}
\begin{itemize}
\item 用户：????????????????????? gamma?
\item 系统：当前配置进度如下：

几何部分已更新结构类型和尺寸参数。  
材料部分需要补充体积与材料的绑定关系。  
源部分已设定粒子类型、能量和位置。  
物理部分需要选择物理列表。  
输出部分需指定格式和保存路径。

请继续补充剩余必要信息。
\item 用户：????? 40 mm??? 80 mm??? G4\_Fe??? gamma 2 MeV??? (0,0,-200) mm??? +z?physics=Shielding?output=csv?
\item 系统：检测到即将覆盖已确认的配置项，请确认是否应用以下修改：

- 几何结构类型：从 single\_box 改为 single\_tubs
- 源能量：从 1.0 改为 2.0
\item 用户：确认
\item 系统：配置已完成。
\end{itemize}
\textbf{内部链路（自然语言）：}
\begin{itemize}
\item 第 1 轮：动作=summarize\_progress，阶段=materials，闭环状态=False。
\item 缺失字段：materials.volume\_material\_map、physics.physics\_list、output.format、output.path。
\item LLM 使用=True；回退=无；后端=llm\_slot\_frame+runtime\_semantic。
\item 第 2 轮：动作=confirm\_overwrite，阶段=source，闭环状态=False。
\item 缺失字段：无。
\item LLM 使用=True；回退=无；后端=llm\_slot\_frame+runtime\_semantic。
\item 第 3 轮：动作=finalize，阶段=source，闭环状态=True。
\item 缺失字段：无。
\item LLM 使用=True；回退=无；后端=llm\_semantic\_frame+runtime\_semantic+confirmed\_pending\_overwrite。
\end{itemize}
\subsection{案例 MTM04（physics）}
\textbf{对话：}
\begin{itemize}
\item 用户：Gamma point source 1 MeV at (0,0,-100) toward +z; please build the rest.
\item 系统：Here's your current setup:

**✅ Source configured:**
- Type: Point source
- Particle: Gamma
- Energy: 1 MeV
- Position: (0, 0, -100)
- Direction: +z

**⏳ Still needed:**
- **Geometry:** Structure type and module dimensions (x, y, z)
- **Materials:** Assignment of materials to volumes
- **Physics:** Physics list selection
- **Output:** File format and save location

Please provide the missing geometry details next.
\item 用户：Set copper box 1 m x 1 m x 1 m, physics FTFP\_BERT, output hdf5.
\item 系统：Your configuration is complete.

**✅ Source:** Gamma point source, 1 MeV, at (0,0,-100) toward +z.
**✅ Geometry:** Copper box, 1 m x 1 m x 1 m.
**✅ Physics:** FTFP\_BERT physics list selected.
**✅ Output:** HDF5 format.
**✅ Materials:** Copper assigned to the main volume.

All necessary components are set. You can proceed with the simulation.
\end{itemize}
\textbf{内部链路（自然语言）：}
\begin{itemize}
\item 第 1 轮：动作=summarize\_progress，阶段=geometry，闭环状态=False。
\item 缺失字段：geometry.structure、geometry.params.module\_x、geometry.params.module\_y、geometry.params.module\_z、materials.volume\_material\_map、physics.physics\_list、output.format、output.path。
\item LLM 使用=True；回退=无；后端=llm\_slot\_frame+runtime\_semantic。
\item 第 2 轮：动作=finalize，阶段=materials，闭环状态=True。
\item 缺失字段：无。
\item LLM 使用=True；回退=无；后端=llm\_slot\_frame+runtime\_semantic。
\end{itemize}
\subsection{案例 MTM05（nuclear\_engineering）}
\textbf{对话：}
\begin{itemize}
\item 用户：?????????????????????
\item 系统：你打算构建哪种几何结构呢？比如一个简单的盒子，还是一个圆柱体？
\item 用户：??????? 200 mm x 200 mm x 200 mm?source.type=point?particle=neutron?energy=14 MeV?position=(0,0,-50) mm?direction=+z?physics=QGSP\_BERT\_HP?output=xml?
\item 系统：好的，当前配置进度如下：

**几何**：已根据你的输入，确定了结构类型和模块尺寸。
**材料**：还需要你指定体积与材料的绑定关系。
**源**：已设置源类型、粒子种类和能量。
**物理**：物理列表及相关选择已同步。
**输出**：输出格式和路径已就绪。

接下来，请告诉我各个几何体积分别对应什么材料。
\end{itemize}
\textbf{内部链路（自然语言）：}
\begin{itemize}
\item 第 1 轮：动作=ask\_clarification，阶段=geometry，闭环状态=False。
\item 缺失字段：geometry.structure、geometry.params.module\_x、geometry.params.module\_y、geometry.params.module\_z、materials.volume\_material\_map、source.type、source.particle、source.energy、source.position、source.direction、physics.physics\_list、output.format、output.path。
\item LLM 使用=True；回退=无；后端=llm\_semantic\_frame+runtime\_semantic。
\item 第 2 轮：动作=summarize\_progress，阶段=materials，闭环状态=False。
\item 缺失字段：materials.volume\_material\_map。
\item LLM 使用=True；回退=无；后端=llm\_slot\_frame+runtime\_semantic。
\end{itemize}
\subsection{案例 MTM06（security\_screening）}
\textbf{对话：}
\begin{itemize}
\item 用户：Need a quick cargo-screening toy setup with forward photons; details can follow.
\item 系统：Here's your current setup status:

**Geometry:** Updated with structure and module dimensions (X, Y).
**Materials:** Still requires volume-to-material binding.
**Source:** Updated with type, particle (photon), and energy.
**Physics:** Still requires a physics list.
**Output:** Still requires format and file path.
\item 用户：Use steel cylinder target radius 60 mm half-length 120 mm, beam gamma 3 MeV from (0,0,-500) mm to +z, physics Shielding, output root.
\item 系统：Your request would overwrite two existing settings. Please confirm this change:

*   **Geometry structure:** `single\_box` → `single\_tubs`
*   **Source energy:** `1.0` → `3.0`
\item 用户：confirm
\item 系统：Your configuration is complete. All required settings have been confirmed, including the output path and format, the Shielding physics list, the material bindings, and the geometry root volume.
\end{itemize}
\textbf{内部链路（自然语言）：}
\begin{itemize}
\item 第 1 轮：动作=summarize\_progress，阶段=materials，闭环状态=False。
\item 缺失字段：materials.volume\_material\_map、physics.physics\_list、output.format、output.path。
\item LLM 使用=True；回退=无；后端=llm\_slot\_frame+runtime\_semantic。
\item 第 2 轮：动作=confirm\_overwrite，阶段=source，闭环状态=False。
\item 缺失字段：无。
\item LLM 使用=True；回退=无；后端=llm\_slot\_frame+runtime\_semantic。
\item 第 3 轮：动作=finalize，阶段=source，闭环状态=True。
\item 缺失字段：无。
\item LLM 使用=True；回退=无；后端=llm\_semantic\_frame+runtime\_semantic+confirmed\_pending\_overwrite。
\end{itemize}
\subsection{案例 MTM07（semiconductor）}
\textbf{对话：}
\begin{itemize}
\item 用户：?????????????????????
\item 系统：你希望构建哪种几何结构呢？比如一个简单的长方体，还是一个圆柱体？
\item 用户：?? 20 mm x 20 mm x 1 mm ?????isotropic ??? 2 MeV ? (0,0,0)?physics=QGSP\_BERT\_HP?output=json?
\item 系统：当前配置进度如下：

- **几何**：已设置结构类型和尺寸参数。
- **材料**：还需补充体积与材料的对应关系。
- **源**：已设置类型、粒子与能量。
- **物理**：已选定物理列表及相关依据。
- **输出**：已确定格式与保存路径。

请继续补充材料绑定信息，完成后即可生成配置文件。
\end{itemize}
\textbf{内部链路（自然语言）：}
\begin{itemize}
\item 第 1 轮：动作=ask\_clarification，阶段=geometry，闭环状态=False。
\item 缺失字段：geometry.structure、geometry.params.module\_x、geometry.params.module\_y、geometry.params.module\_z、materials.volume\_material\_map、source.type、source.particle、source.energy、source.position、source.direction、physics.physics\_list、output.format、output.path。
\item LLM 使用=True；回退=无；后端=llm\_semantic\_frame+runtime\_semantic。
\item 第 2 轮：动作=summarize\_progress，阶段=materials，闭环状态=False。
\item 缺失字段：materials.volume\_material\_map。
\item LLM 使用=True；回退=无；后端=llm\_slot\_frame+runtime\_semantic。
\end{itemize}
\subsection{案例 MTM08（education）}
\textbf{对话：}
\begin{itemize}
\item 用户：For classroom demo, I want gamma attenuation in copper; you can ask only essential missing parts.
\item 系统：Here's your current setup:

**Geometry:** Still needed — the structure type and module X dimension.
**Materials:** Copper is selected. Still needed — which volume uses this material.
**Source:** Gamma particle is selected. Still needed — the source type and its energy.
**Physics:** Physics list is configured.
**Output:** Still needed — the format and file path for results.

Please provide the missing geometry, source, and output details.
\item 用户：Okay, copper box 100 mm x 100 mm x 100 mm, point gamma 1 MeV at (0,0,-200) mm, direction +z, physics FTFP\_BERT, output csv.
\item 系统：An existing confirmed value would be overwritten. Please confirm this change: physics list from "Shielding" to "FTFP\_BERT".
\item 用户：confirm
\item 系统：Your configuration is now complete.
\end{itemize}
\textbf{内部链路（自然语言）：}
\begin{itemize}
\item 第 1 轮：动作=summarize\_progress，阶段=geometry，闭环状态=False。
\item 缺失字段：geometry.structure、geometry.params.module\_x、geometry.params.module\_y、geometry.params.module\_z、materials.volume\_material\_map、source.type、source.energy、source.position、source.direction、output.format、output.path。
\item LLM 使用=True；回退=无；后端=llm\_slot\_frame+runtime\_semantic。
\item 第 2 轮：动作=confirm\_overwrite，阶段=materials，闭环状态=False。
\item 缺失字段：无。
\item LLM 使用=True；回退=无；后端=llm\_slot\_frame+runtime\_semantic。
\item 第 3 轮：动作=finalize，阶段=source，闭环状态=True。
\item 缺失字段：无。
\item LLM 使用=True；回退=无；后端=llm\_semantic\_frame+runtime\_semantic+confirmed\_pending\_overwrite。
\end{itemize}
\subsection{案例 MTM09（environmental\_radiation）}
\textbf{对话：}
\begin{itemize}
\item 用户：???????????????????????
\item 系统：你希望构建哪种几何结构呢？比如一个简单的盒子，还是一个圆柱体？另外，这个结构在X方向的尺寸是多少？
\item 用户：????? 500 mm x 500 mm x 500 mm?gamma ?? 0.5 MeV??? (0,0,-300) mm??? +z?physics=QBBC?output=root?
\item 系统：好的，当前配置进度如下：

**几何**：已设置好结构类型和X、Y方向的尺寸。
**材料**：还需要补充体积与材料的绑定关系。
**源**：已设置好类型、粒子种类和能量。
**物理**：物理列表及相关选择已同步。
**输出**：输出格式和文件路径已设置。

接下来需要您确认或补充材料部分的绑定信息。
\end{itemize}
\textbf{内部链路（自然语言）：}
\begin{itemize}
\item 第 1 轮：动作=ask\_clarification，阶段=geometry，闭环状态=False。
\item 缺失字段：geometry.structure、geometry.params.module\_x、geometry.params.module\_y、geometry.params.module\_z、materials.volume\_material\_map、source.type、source.particle、source.energy、source.position、source.direction、physics.physics\_list、output.format、output.path。
\item LLM 使用=True；回退=无；后端=llm\_semantic\_frame+runtime\_semantic。
\item 第 2 轮：动作=summarize\_progress，阶段=materials，闭环状态=False。
\item 缺失字段：materials.volume\_material\_map。
\item LLM 使用=True；回退=无；后端=llm\_slot\_frame+runtime\_semantic。
\end{itemize}
\subsection{案例 MTM10（high\_energy\_physics）}
\textbf{对话：}
\begin{itemize}
\item 用户：Set geometry as a sphere target; I have not decided radius or source details yet.
\item 系统：Here's a summary of your Geant4 configuration so far:

**Geometry:**  
You've set the structure to a sphere target and defined the root volume. We still need the sphere's maximum radius (child\_rmax).

**Materials:**  
We need to assign materials to your volumes.

**Source:**  
The particle source details are still needed: source type, particle type, and energy.

**Physics \& Output:**  
A physics list, output format, and output path are also required to complete the setup.
\item 用户：Sphere target radius 80 mm, lead material, electron point source 5 MeV at (0,0,-200) mm toward +z, physics FTFP\_BERT, output hdf5.
\item 系统：Your Geant4 configuration is complete.

You have set a sphere target with an 80 mm radius, assigned lead as the material, and defined an electron point source with 5 MeV energy. The physics list is FTFP\_BERT, and the output format is HDF5.
\end{itemize}
\textbf{内部链路（自然语言）：}
\begin{itemize}
\item 第 1 轮：动作=summarize\_progress，阶段=geometry，闭环状态=False。
\item 缺失字段：geometry.params.child\_rmax、materials.volume\_material\_map、source.type、source.particle、source.energy、source.position、source.direction、physics.physics\_list、output.format、output.path。
\item LLM 使用=True；回退=无；后端=llm\_slot\_frame+runtime\_semantic。
\item 第 2 轮：动作=finalize，阶段=materials，闭环状态=True。
\item 缺失字段：无。
\item LLM 使用=True；回退=无；后端=llm\_slot\_frame+runtime\_semantic。
\end{itemize}
\end{document}